%% kapa: exam=1 round=33 year=2022 period=1 form=A subject=civil_law subject_ko=민법 seq=20 number=20
%% source: Q-Net 감정평가사 기출문제 (cst00309) artlSeq=5209997 file=raw/1008/5209997/2022년도 제33회 감정평가사 1차시험 1교시(민법,경제학,부동산학).pdf
%% conversion: kordoc 4.24.51 -> tools/gen1.py (auto); review=proofread-sonnet; flags=none
\begin{question}{20}
소멸시효의 중단과 정지에 관한 설명으로 옳지 않은 것은?
\begin{choices}
  \item 시효의 중단은 원칙적으로 당사자 및 그 승계인 간에만 효력이 있다.
  \item 파산절차참가는 채권자가 이를 취소하거나 그 청구가 각하된 때에는 시효중단의 효력이 없다.
  \item 부재자재산관리인은 법원의 허가 없이 부재자를 대리하여 상대방의 채권의 소멸시효를 중단시키는 채무의 승인을 할 수 없다.
  \item 천재 기타 사변으로 인하여 소멸시효를 중단할 수 없을 때에는 그 사유가 종료한 때로부터 1월내에는 시효가 완성하지 아니한다.
  \item 부부 중 한쪽이 다른 쪽에 대하여 가지는 권리는 혼인관계가 종료된 때부터 6개월 내에는 소멸시효가 완성되지 아니한다.
\end{choices}
\end{question}
\answer{3}
